% status: ZERO
% Chapter 10 of the second edition. Plan: docs/STRUCTURE.md. Rules: AUTHORING.md.
% Measurements: research/measurements/2026-09-23-m1-part2.md
\chapter{Quantizing Weights: INT8, INT4, GPTQ, AWQ and K-Quants}\label{ch:weight-quantization}

\begin{ChapterIntent}
Batch-one decode reads every weight once per token, so storing weights in
four bits instead of sixteen makes decode almost four times faster and lets a
model four times larger fit. This chapter builds weight quantization from
scales and groups upward: round-to-nearest and its failure, GPTQ's
error-correcting updates, AWQ's protection of salient channels,
weight--activation schemes such as SmoothQuant, and the K-quant formats of
llama.cpp. It then puts dequantization inside a kernel, where bandwidth savings
can be lost to arithmetic, and shows how to measure what quantization costs in
accuracy.
\end{ChapterIntent}

\OpeningSection{Fewer bits, more tokens}

Take one model, Llama~3.2~3B, and store it six ways, from llama.cpp's 8-bit
format down to its 2-bit K-quant. On the M1 of \Chref{roofline}, the files
range from 3.41~GB to 1.36~GB, and decode speed runs the other way: from
\textbf{15.0 to 28.8 tokens per second} as the file shrinks
(measured; record \texttt{2026-09-23-m1-part2};
Figure~\ref{fig:ch10-quant-sweep}). Multiply each file's size by its tokens
per second and the product --- the bandwidth it effectively reads at --- stays
between 39 and 51~GB/s. Decode speed is, to first order, the inverse of the
bytes in the file. (A first run of the same sweep, while the machine was busy
with other work, measured every format about 20 to 30 percent slower; the
record keeps both runs, and the numbers here are from the second.)

\EngineFigure{ch10-quant-sweep}{Decode speed against file size for six
llama.cpp formats of Llama~3.2~3B on the M1 (measured, second run). Each curve
is a constant effective bandwidth: points on one curve read their bytes equally
fast. The 8-bit file sits on the 51~GB/s curve, the 3- and 2-bit files on the
39~GB/s curve.}{fig:ch10-quant-sweep}

Prefill does not care. The same six files process a 512-token prompt at
246--281 tokens per second, with no trend in size: prefill is compute-bound
(\Chref{prefill-decode}), and fewer bytes per weight do not reduce the
arithmetic.

Two caveats frame the rest of the chapter. First, these files were made by
re-quantizing the 4-bit file, so they measure speed only; their accuracy is
not representative of the formats. llama.cpp's own quantizer publishes the
accuracy side: on Llama-3-8B its table lists perplexity increases of 0.02 for
6-bit, 0.06 for 5-bit, 0.66 for 3-bit and 3.52 for 2-bit K-quants (the
quantizer's built-in help, build \texttt{d00685831}; the maintainers'
measurement, not reproduced here). Second, the smallest formats are not
proportionally faster. The 8-bit file is read at 51~GB/s, about what a plain
streaming kernel reached in \Chref{roofline}; the 4- to 6-bit K-quants at
46--48; the 3- and 2-bit files at 39. Q3\_K\_M is 17\% smaller than Q4\_K\_M
and only 3\% faster. Unpacking the weights has started to cost arithmetic that
the bandwidth no longer hides. That trade --- bytes saved against work added
and accuracy lost --- is this chapter's subject.

\NapkinSection{Napkin math: bits per weight to tokens per second}

A weight-quantized model reads $P\,b_w/8$ bytes per decode step, where $b_w$
is the average bits per weight \emph{including} scales. So
\begin{equation}
r_{\max} = \frac{8\beta}{P\,b_w}\quad[\text{tokens/s}],
\label{eq:quant-ceiling}
\end{equation}
and halving $b_w$ doubles the ceiling. The averages are larger than format
names suggest. llama.cpp's ``Q4\_K\_M'' Llama~3.2~3B file averages
\textbf{5.0 bits per parameter}: its K-quant blocks cost 4.5 bits per weight
including scales \cite{kawrakow2023kquants}, and the mix keeps the vocabulary
table and, in half the layers, the value and FFN-down projections in the
6.56-bit Q6\_K format (read from the file; record
\texttt{2026-09-23-gguf-tensor-bytes}). For a 70-billion-parameter model on an
H100, equation~\eqref{eq:quant-ceiling} gives 24 tokens per second at 16 bits,
48 at 8 and about 85 at a real 4.5-bit average --- if it fits: only the last
two do, in 80~GB.

\section{What a quantized weight is}

A quantized weight is an integer code $q$ with $b$ bits and a rule for turning
it back into a number. The usual rule is affine:
\begin{equation}
w \approx s\,(q - z),\qquad q\in\{0,\dots,2^{b}-1\},
\label{eq:affine-quant}
\end{equation}
with a scale $s$ and a zero point $z$ shared by a group of weights. Quantizing
chooses $s$, $z$ and each $q$; dequantizing computes $s(q-z)$ inside the
kernel, just before the multiply. A symmetric scheme fixes $z$ at the middle of
the range and stores only $s$.

\section{Scales, zero points and groups}

One scale for a whole matrix wastes the grid on its few largest values. One
scale per group of $g$ consecutive weights adapts to local ranges at a cost of
$(b_s+b_z)/g$ extra bits per weight --- 0.25 bits for a 16-bit scale and a
16-bit zero shared by 128 weights. llama.cpp's K-quants go a level further:
super-blocks of 256 weights, split into sub-blocks with their own scales, which
are themselves quantized to 6 bits against a super-block scale
\cite{kawrakow2023kquants}. Q4\_K spends 4.5 bits per weight this way; Q6\_K,
6.5625; Q2\_K, 2.5625.

\section{Round-to-nearest and why it fails}

The obvious quantizer rounds each weight to its nearest grid point. It is fast
and needs no data, and at 8 bits it is nearly lossless. At 4 bits and below it
fails, for a reason that is not about weights at all: the error that matters is
in the layer's \emph{output}, $\mathbf{W}\mathbf{x}$, and a few input channels
of $\mathbf{x}$ carry much larger values than the rest. Small rounding errors
in the weights that multiply those channels are amplified; large errors
elsewhere barely matter. Every method below uses data about the activations to
decide where the error goes.

\section{GPTQ: second-order error correction}

GPTQ quantizes a weight matrix one column at a time and, after each column,
adjusts the columns not yet quantized to compensate for the error just made,
using second-order information estimated from a small set of calibration
activations \cite{frantar2023gptq}. It is one-shot --- no retraining --- and
its authors quantized 175-billion-parameter models to 3--4 bits in about four
GPU hours with little loss.

\section{AWQ: protect the salient channels}

AWQ observes that protecting about 1\% of weights --- those that multiply the
largest-magnitude activation channels --- removes most of the error, and that
it can protect them without storing them separately: scale those input
channels up in the weights and down in the activations, an equivalent
transformation, before quantizing \cite{lin2024awq}. It needs no gradient or
reconstruction, only activation statistics; its authors ran the resulting
4-bit models, up to Llama-2-70B, on desktop and mobile GPUs with their own
runtime.

\section{Weights and activations: SmoothQuant and W8A8}

Quantizing activations too lets the matrix product itself run in low-precision
integer arithmetic rather than dequantizing to 16 bits first --- a compute
saving for prefill, where compute binds. Activations are harder to quantize
than weights because of those outlier channels. SmoothQuant migrates the
difficulty into the weights with the same kind of equivalent scaling, enabling
8-bit weights and activations \cite{xiao2023smoothquant}. At 4 bits,
rotation-based methods multiply activations and weights by orthogonal matrices
that spread outliers evenly across channels without changing the layer's
output: QuaRot uses Hadamard rotations and quantizes weights, activations and
KV cache all to 4 bits \cite{ashkboos2024quarot}; SpinQuant learns the
rotations \cite{liu2025spinquant}.

\section{K-quants and mixed-precision files}

A quantized checkpoint need not use one format. llama.cpp's mixes assign
formats per tensor type and layer, spending bits where the model is most
sensitive: in the Q4\_K\_M file above, 6-bit K-quants for the vocabulary table
and for value and FFN-down projections in half of the layers, 4.5-bit K-quants
for everything else, and 32-bit floats for norms. The name describes the
majority format, not the average.

\section{Dequantization inside the kernel}

A weight-only quantized matrix product loads packed codes, unpacks them,
applies scales, and multiplies in 16- or 32-bit arithmetic --- extra work per
weight that the bandwidth saving must pay for. At batch one there is spare
arithmetic, and the kernel stays bandwidth-bound; the opening measurement's
46--51~GB/s for the 4- to 8-bit files is that regime, and the fall to
39~GB/s at 2 and 3 bits is the unpacking beginning to show. As the batch grows, the unpacking is repeated
work per weight that batching does not amortize the same way, and a
quantized kernel can become compute-bound before a 16-bit one would.
Specialized kernels push that point out: Marlin, an FP16$\times$INT4 kernel,
stays close to the ideal 4$\times$ speedup up to batches of 16--32
\cite{frantar2024marlin}, and vLLM's Machete kernel does the same job with
Hopper's asynchronous copies and warpgroup MMA (vLLM pull request 7174). The
fastest CPU kernels avoid floating point altogether and compute dot products
directly on the integer codes, applying scales once at the end.

\section{Measuring what you lost}

Quantization's cost is a change in the model's output distribution, and the
honest way to measure it is to compare distributions. Perplexity on a text
corpus is a coarse average; the Kullback--Leibler divergence between the
quantized and unquantized models' next-token distributions, token by token,
shows where and how often they disagree; task evaluations show whether the
disagreements matter. \Chref{fast-wrong-answers} builds the harness. A
quantized model whose perplexity barely moves can still change the answers to
the questions a particular application asks.

\LedgerSection

Weight-only quantization from 16 bits to llama.cpp's Q4\_K\_M mix, 5.0 bits per weight on average:

\begin{ConstraintLedger}
Memory bandwidth & buys & About 3.2$\times$ fewer bytes per decode step than 16-bit weights (derived); on the M1, going from 8-bit to Q4\_K\_M cut the bytes 1.70$\times$ and raised decode speed 1.52$\times$ (measured). \\
Memory capacity & buys & The same factor in resident weight bytes. \\
Compute & spends & Unpacking and scaling per weight; visible at 2--3 bits and at larger batches. \\
Correctness & risks & Perplexity and task accuracy fall as bits fall; measure per model (llama.cpp's own table: +0.66 perplexity at 3 bits on Llama-3-8B). \\
\end{ConstraintLedger}

\LabSection{quantize a matrix three ways}

\LabIntent{In \texttt{code/labs/ch10-weight-quant/}, a Rust program that takes
a real weight matrix (dequantized from a GGUF file) and real activations from
the mini-engine, quantizes the weights by per-tensor round-to-nearest,
per-group round-to-nearest, and an AWQ-style activation-aware scaling, and
reports the error in the layer's \emph{output} rather than in the weights. A
second part times a dequantizing matrix--vector product against the 32-bit
one, checking both against a float64 oracle, and finds the batch size at which
the quantized kernel stops winning.}

\HappenedSection{a SIMD matrix--vector kernel slower than scalar}

Hermon's expert-streaming work needed a matrix--vector product over K-quant
weights, and its first version was scalar. The obvious vectorization ---
unpack each 16- or 32-weight sub-block to floats and call an already
vectorized dot product --- made it \textbf{slower}: 0.62 tokens per second
against the scalar 0.72 on the same trace. Every short dot product ended in a
horizontal reduction across the vector lanes, and on 16 elements the reduction
cost more than the multiply--adds it summed. Folding each sub-block's scale into
the unpacking let a whole 256-weight super-block be reduced once, eight times
fewer reductions, and the result was about \textbf{1.3$\times$} faster than
scalar on an x86 machine with AVX2 --- far from the 8--19$\times$ the same
team had measured vectorizing other kernels (Hermon
\texttt{docs/STORAGE\_ARCHITECTURE.md}, section 5.8, commit \texttt{a7c0597}).
Its authors located the remaining multiple precisely: the fast implementations
never convert the codes to floating point, computing integer dot products and
scaling once at the end. Correctness held throughout: the vectorized path
reproduced a NumPy reference on real Qwen3 weights exactly. The kernel is
callable through the kernel crate's public interface, but no request path
calls it at commit \texttt{2a3fd52}, so LIBRARY.

\FrontierSection{2026-09-24}

Post-training weight quantization is becoming one option among several.
Models increasingly arrive already quantized or trained for a format
(\Chref{low-precision-float}); rotation-based methods push weights,
activations and cache together to 4 bits \cite{ashkboos2024quarot,liu2025spinquant};
and serving engines carry specialized mixed-precision kernels per GPU
generation \cite{frantar2024marlin}. At the low end, llama.cpp's i-quants
replace uniform grids with fixed code books: an index selects a pattern of
eight (IQ2) or four (IQ3) weight magnitudes from a 256-entry table, and the IQ4
formats use sixteen non-uniform levels. Scales included, they run from 1.5625
bits per weight (IQ1\_S) through 2.0625 (IQ2\_XXS) and 3.0625 (IQ3\_XXS) to
4.25 (IQ4\_XS) (block definitions in llama.cpp
\texttt{ggml/src/ggml-common.h} at commit \texttt{389ff61d}).

\ExercisesSection

\begin{enumerate}
\item \emph{Derivation.} Compute the average bits per weight of a file that
stores 90\% of its weights in Q4\_K and 10\% in Q6\_K. What decode speed does
equation~\eqref{eq:quant-ceiling} predict on your machine?
\item \emph{Prediction.} Before running the lab, predict the batch size at
which a dequantizing kernel stops beating the 16-bit one on your CPU.
\item \emph{Implementation.} Add group-wise zero points to the lab's
round-to-nearest quantizer and measure whether they help at 4 and 3 bits.
\item \emph{Falsification.} Construct a weight matrix and activation
distribution for which round-to-nearest beats an activation-aware scaling.
\end{enumerate}
